\section{Confluence of action, area, information, and gravitation}
\label{sec:confluencia}

The area scale in this section is the same one determined by the
longitudinal and radial composition in \ref{g26:sec:rigidez-radial-en}:
\begin{equation}
 \ell_P^2=L^2=\frac{G_a\hbar_a}{c_{\rm int}^3},\qquad
 E_{P,a}=\frac{\hbar_a c_{\rm int}}{L}
        =k_B\Theta_{{\rm clk},a}\log3.
 \label{eq:ii-g26-longitud-energia-area}
\end{equation}
In the return section, \(E_{P,\rm ret}=E_L\). The variable horizon
radius \(R\) retains its own meaning; condition \(R=L\) specifies
the particular cell studied below. The increment
\(\delta\mathcal A=\gamma_{\HMT}a_0L^2\) transmits that same area
unit to incidence, retaining action, the thermal transducer, and the
occupations of the reduced state.

\subsection{Energy and entropy of the half-action cell}

Fix the same oriented action section and its corresponding
gravitational and thermal sections. The preceding reciprocity allows the same cell to be expressed through
action, duration, and length. With
\(\ell_P^2=\hbar G/c^3\), \(t_P=\ell_P/c\), and
\(E_P=\hbar/t_P\), one has
\[
 \frac{c^4\ell_P}{G}=\frac{\hbar}{t_P}=E_P.
\]
The equality follows by substituting the area scale already constructed. In the
horizon realization defined by \eqref{eq:ii-horizonte-normalizado}, the quantities
are normalized as
\begin{equation}
 E_\Lambda(R)=\frac{c^4R}{2G},\qquad
 S_H(R)=\frac{k_B\pi R^2}{\ell_P^2},\qquad
 T_H(R)=\frac{\hbar c}{2\pi k_BR}.
 \label{eq:ii-horizonte-normalizado}
\end{equation}
The temperature corresponds to the cosmological-horizon normalization
of this realization. The factor \(2\pi\) and the definition of the radius
are part of the set of equations used.

\begin{theorem}[Identity of the gravitational half-action cell]
In realization \eqref{eq:ii-horizonte-normalizado},
\begin{align}
 T_H(R)S_H(R)&=E_\Lambda(R),\\
 T_H(\ell_P)S_H(\ell_P)&=\frac{E_P}{2}
 =\frac{k_B\Theta_{\rm clk}\log3}{2},\\
 E_\Lambda(\ell_P)t_P&=\frac{\hbar}{2},\qquad
 E_\Lambda(\ell_P)T_{108}=\frac h2.
\end{align}
\end{theorem}
\begin{proof}
Multiplying temperature by entropy cancels the thermal
normalization and gives \(\hbar cR/(2\ell_P^2)=c^4R/(2G)\).
At \(R=\ell_P\), the expression is \(\hbar c/(2\ell_P)
=\hbar/(2t_P)=E_P/2\). The thermal identity of the cycle gives
the second member. Finally, \(T_{108}=2\pi t_P\) and
\(h=2\pi\hbar\) produce the two action relations.
\end{proof}

The theorem brings together quantities that previously belonged to different
descriptions. The cell energy, the temperature--entropy product,
tritic information, and action integrated over its duration are
compatible readings under the declared normalization. The gravitational
constant enters the area scale; action fixes the energy
of the period; Boltzmann transports dimensionless information to the
entropy line. The proof shows how these functions are
joined in a single identity.

\subsection{Conversion between decision energy and area increment}
\label{sec:ii-conversion-informacion-area}

Energy per tritic decision and the elementary area increment are
obtained through different readers in the preceding development:
\begin{equation}
 E_P=k_B\Theta_{\rm clk}\log3=\frac{\hbar}{t_P},
 \qquad
 \delta\mathcal A=\gamma_{\HMT}a_0\ell_P^2.
 \label{eq:ii-confluencia-cuanto-energia-area}
\end{equation}
The same oriented section and the positive branch
\(\hbar,G,\gamma_{\HMT},a_0,t_P>0\) are maintained. The first quantity
belongs to the energy line and the second to the area line.
Their quotient defines the conversion coefficient
\begin{equation}
 \mathcal T_{I/A}
 :=\frac{E_P}{\delta\mathcal A}
 =\frac{k_B\Theta_{\rm clk}\log3}
        {\gamma_{\HMT}a_0\ell_P^2},
 \qquad [\mathcal T_{I/A}]=[\text{energy}]/[\text{area}].
 \label{eq:ii-confluencia-tension}
\end{equation}
This definition makes explicit the normalization required to transport
a decision energy to a geometric resolution. All its factors
are specified before the quotient is formed.

\begin{proposition}[Gravitational expression of the conversion coefficient]
\label{prop:ii-confluencia-tension}
In the same dimensional realization,
\begin{equation}
 \mathcal T_{I/A}
 =\frac{c^3}{\gamma_{\HMT}a_0G t_P}
 =\frac{c^4}{\gamma_{\HMT}a_0G\ell_P}.
 \label{eq:ii-confluencia-tension-gravedad}
\end{equation}
\end{proposition}
\begin{proof}
Substituting \(E_P=\hbar/t_P\) and \(\ell_P^2=\hbar G/c^3\)
into \eqref{eq:ii-confluencia-tension} cancels the common action
section and gives the first member. The identity \(\ell_P=ct_P\)
produces the second. The dependence of \(t_P\) and
\(\ell_P\) on the previously constructed realization is preserved.
\end{proof}

In the minimal cut of \eqref{eq:ii-area-corte-informacion}, the
realization of one uniform decision per link gives
\begin{equation}
 S_{\rm corte}^{\rm fis}=81k_B\log3,
 \qquad
 E_{\rm corte}^{(1)}=81E_P
 =\Theta_{\rm clk}S_{\rm corte}^{\rm fis}.
 \label{eq:ii-confluencia-corte-energia}
\end{equation}
The first equality uses the proved capacity and the uniform product
state; the second assigns the decision energy of the same cycle to
each of its eighty-one links. Multiplying the thermal identity
by eighty-one proves the equation of state of this realization.
The thirty-six-link sector state preserves its own
partition function, entropy, and area variations in
\eqref{eq:ii-area-ley-finita}.

The composition identifies the complementary functions of the
constants: action fixes energy per duration; the angular character
and Barbero--Immirzi fix the sector resolution of area; Boltzmann
dimensionally realizes information; the product \(G\hbar\)
determines the area unit. The coefficient \(\mathcal T_{I/A}\)
preserves these dependencies by expressing their relation in units of
energy per area.

\subsection{Physical meaning of area confluence}

For a horizon of radius \(R\), the area is \(A_H=4\pi R^2\).
The same entropy relation admits two expressions
\[
 \frac{S_H}{k_B}=\frac{A_H}{4\ell_P^2}
 =\frac{\pi R^2c^3}{G\hbar}.
\]
The first quotient compares the macroscopic area with the construction's
area unit. The second makes explicit the composition of that unit
through action and gravitation. Transporting the sections so
that \(G_a\hbar_a\) remains fixed makes the geometric reading preserve
this normalization. The dimensionless entropy of the same horizon remains
invariant; its dimensional expressions preserve the bases they
use.

The Barbero--Immirzi parameter acts at a different resolution
of the same development: oriented incidence organizes the elementary
area spectrum, and the reduced state describes the information
accessible at the cut. The areal and volumetric sheets provide the
geometric transport and history measure described in
\cref{sec:ii-hojas-av}. Comparing a microscopic spectrum
with a horizon entropy retains its realization map,
state, and scale as data. These specifications allow the
physical relation to be interpreted at each level and avoid replacing a correspondence
of observables with an indiscriminate identification of their spaces.

The theoretical unity of the argument resides in complementary functions:

The area unit combines the product $G\hbar$; oriented incidence
organizes its elementary resolution through $\gamma_{\HMT}a_0$; the
state and bipartition determine the information accessible to the
subsystem. Response equations
\eqref{eq:ii-respuesta-entropica}--\eqref{eq:ii-respuesta-areal}
make this articulation explicit in the sector collection. The inversion
\eqref{eq:ii-inversion-entropia} preserves entropy and transforms
area according to~\eqref{eq:ii-area-complementaria}: the two readings
distinguish different properties of the same configuration.

In this articulation,
the dynamics produces and preserves history; the reader determines what
information is accessible; action relates evolution and energy;
Boltzmann realizes its thermal expression; gravitation enters
the geometric scale by which area is measured. The identity
\(T_HS_H=E_\Lambda\) shows an exact confluence of these functions
in the declared realization. Physical theory is thus formulated from
mathematical objects and their effective transformations.

\subsection{Spectral entropy and preservation under transport}

The entropy of a reduced state depends on its spectrum and on the
bipartition defining the reduction. To specify continuation
toward duality, let \(\mathcal H_A,\mathcal H_B\) be the two factors
and \(\rho_{AB}\) a state with finite entropies. A factorized
transport \(U_A\otimes U_B\) produces
\[
 \rho'_A=U_A\rho_AU_A^*,\qquad S(\rho'_A)=S(\rho_A).
\]
The first identity follows from the definition of partial trace, and the
second from spectrum invariance. For a general unitary
transport, the equivalent formulation also transports
the inclusion of the subsystem's observable algebra. This qualification
identifies the geometric datum that must accompany a claim
of entanglement preservation.

The reciprocity of the ellipse and momentum--winding exchange
preserve the quantities indicated in their respective theorems. Transport
of the bipartition determines how to use these identities
in an entropy reading. Modular continuation thus
maintains the article's informational content and keeps distinct
the energy spectrum, action orientation, and algebraic reduction
defining entanglement.

\subsection{Effective multiplicity and hyperbolic realization}

In the preceding cell, \(S_H(\ell_P)/k_B=\pi\). The effective
multiplicity associated with that entropy is
\[
 \Omega_{\rm eff}=\exp(S_H/k_B)=e^\pi.
\]
If \(H\) is the hyperbolic involution of TRIT, with eigenvalues
\(+1,-1\), the exponential \(\exp(\pi H)\) has spectrum
\(\{e^\pi,e^{-\pi}\}\). Agreement of the expansive eigenvalue
with the effective multiplicity locates a specific spectral relation.
Its microscopic meaning retains as additional data the map between
the horizon space and the hyperbolic representation. This
localization allows the correspondence to be developed without confusing
the spectral value with its entire realization space.
